%% -------------------------------------------------------
%% Anhang B: Dokumentation der KI-Nutzung
%% -------------------------------------------------------
\chapter{Dokumentation der KI-Nutzung}
\label{app:ki}

Gemäß den Richtlinien der HTL Leonding müssen alle im Rahmen dieser Diplomarbeit
verwendeten KI-generierten Inhalte dokumentiert werden. Die folgende Tabelle listet
alle Einsätze auf, bei denen ein KI-Tool zur Unterstützung herangezogen wurde.

\begin{longtable}{|p{0.12\textwidth}|p{0.25\textwidth}|p{0.35\textwidth}|p{0.18\textwidth}|}
  \hline
  \textbf{Datum} & \textbf{KI-Tool} & \textbf{Verwendungszweck / Prompt (Zusammenfassung)} &
  \textbf{Kapitel / Datei} \\
  \hline
  \endfirsthead
  \hline
  \textbf{Datum} & \textbf{KI-Tool} & \textbf{Verwendungszweck} & \textbf{Kapitel} \\
  \hline
  \endhead
  % Einträge hier ergänzen – Beispiel:
  [Datum] & Claude (Anthropic) &
  Strukturvorschlag für Kapitel Systemarchitektur; Überarbeitung auf Basis eigener
  Inhalte und Projektcode &
  Kap.~\ref{chap:architektur} \\
  \hline
  [Datum] & Claude (Anthropic) &
  Vorlage für JWT-Authentifizierungslogik; angepasst an eigene Modelle und Konfiguration &
  Kap.~\ref{chap:backend} \\
  \hline
  [Datum] & Claude (Anthropic) &
  Entwurf der SwiftUI-Listings; geprüft und an eigene Projektstruktur angepasst &
  Kap.~\ref{chap:ios} \\
  \hline
  \multicolumn{4}{|p{0.9\textwidth}|}{
    \textit{Alle KI-generierten Inhalte wurden kritisch geprüft, inhaltlich angepasst
    und auf ihre Korrektheit hin verifiziert. Die endgültige Verantwortung für den Inhalt
    liegt beim Autor dieser Arbeit.}
  } \\
  \hline
\end{longtable}
